\scanpage{202}\mspage{151}
soit iso. à la catégorie ponctuelle, i.e $\mathrm{Hom}_X(t_0,t_0)=\{\mathrm{id}_{\unsure{t|t_0}}\}$. Supposons que pour tout $t\in T\setminus\{t_0\}$, $U_{/t}$ soit $W$-asphérique. Pour que l'on ait $i\in W$, il f et s. que $i$ soit asphérique, i.e que $U_{/t_0}$ soit également $W$-asphérique.

Il faut prouver que $i\in W$ implique que $U_{/t_0}$ est $W$-asphérique. Considérons l'inclusion

\[
i^{\prime}_0:U\longrightarrow \deleted{\unsure{X}}U'\overset{\mathrm{déf}}{=}X\setminus\{t_0\},
\]

l'hypothèse sur les $U_{/t}$ \inserted{pour $t\in T-\{t_0\}$} signifie que $i'$ est $W$-asphérique, donc $i'$ est dans $W$, et comme \deleted{\unsure{le}} le composé $i:U\xrightarrow{\unsure{i'_0|i_0}}U'\xrightarrow{i'}X$ est dans $W$, cela implique que $i':U'\mathbin{\unsure{-|\to}}X$ est dans $W$. Nous sommes maintenant dans la situation du cor. 4, qui implique que $i'$ est $W$-asphérique. Donc $i=i'i'_0$ est $W$-asphérique, comme composé de deux morphismes $W$-asphériques, qed.

\uline{Démonstration du \deleted{prop. \unsure{15|11}} \inserted{th. 4}} \textcircled{a} Critère de lissité de

\diagram{\(\varphi_Q:\unsure{\mathscr X|\mathcal X}=\int \begin{tikzcd}[ampersand replacement=\&,row sep=0pt,column sep=1.2em,cramped] \& X_1 \\ X_0\arrow[ur,dash]\arrow[dr] \& \\ \& X_2 \end{tikzcd}\longrightarrow Y.\)}

Comme $\{1\},\{\unsure{1|2}\}$ sont des \deleted{\unsure{sous-cat.}} ouverts de

\diagram{\(\Phi={}\begin{tikzcd}[ampersand replacement=\&,row sep=0pt,column sep=1.4em]
\& 1\arrow[dl]\\0\&\\\&2\arrow[ul]\end{tikzcd},\)}

\deleted{\unsure{$X_1$}}, $i_1:X_1\to\mathscr X$ et $i_2:X_2\to\mathscr X$ sont des imm. ouvertes. Comme $\varphi_Q\circ i_{\unsure{\nu|v|j}}=f_\nu$,
% archive sheet number: 202
